\documentclass[11pt,a4paper]{article}
\usepackage[utf8]{inputenc}
\usepackage[T1]{fontenc}
\usepackage[spanish,es-noindentfirst,es-tabla]{babel}
\usepackage{lmodern}
\usepackage{amsmath,amssymb}
\usepackage{booktabs}
\usepackage{tabularx}
\usepackage{graphicx}
\usepackage{listings}
\usepackage{xcolor}
\usepackage[hidelinks]{hyperref}
\usepackage{cite}
\usepackage{float}
\usepackage{geometry}
\geometry{a4paper, margin=2.2cm, headheight=14pt}
\usepackage{array}
\usepackage{enumitem}
\usepackage{caption}
\usepackage{fancyhdr}

\pagestyle{fancy}
\fancyhf{}
\fancyhead[L]{\small\color{gray} UTEQ -- Aplicaciones Distribuidas [20701]}
\fancyhead[R]{\small\color{gray} PFC E4: Microservicio Soporte y Observabilidad Distribuida}
\fancyfoot[C]{\thepage}
\renewcommand{\headrulewidth}{0.4pt}
\renewcommand{\footrulewidth}{0.4pt}

\definecolor{codegreen}{rgb}{0.15,0.55,0.15}
\definecolor{codegray}{rgb}{0.45,0.45,0.45}
\definecolor{backcolour}{rgb}{0.97,0.98,0.99}
\definecolor{codeblue}{rgb}{0.1,0.25,0.7}
\definecolor{codered}{rgb}{0.75,0.1,0.1}

\lstdefinelanguage{json}{
  basicstyle=\ttfamily\scriptsize,
  showstringspaces=false,
  breaklines=true,
  stringstyle=\color{codered},
  keywordstyle=\color{codeblue}\bfseries,
  commentstyle=\color{codegreen}\itshape,
  morestring=[b]",
  morestring=[d]'
}

\lstdefinelanguage{properties}{
  keywords={management,metrics,distribution,percentiles-histogram,http,server,requests},
  keywordstyle=\color{codeblue}\bfseries,
  basicstyle=\ttfamily\scriptsize,
  sensitive=false,
  comment=[l]{\#},
  commentstyle=\color{codegreen}\itshape,
  stringstyle=\color{codered}
}

\lstdefinestyle{codigo}{
  backgroundcolor=\color{backcolour},
  commentstyle=\color{codegreen}\itshape,
  keywordstyle=\color{codeblue}\bfseries,
  numberstyle=\tiny\color{codegray},
  stringstyle=\color{codered},
  basicstyle=\ttfamily\scriptsize,
  breakatwhitespace=false,
  breaklines=true,
  captionpos=b,
  keepspaces=true,
  numbers=left,
  numbersep=5pt,
  showspaces=false,
  showstringspaces=false,
  showtabs=false,
  tabsize=2,
  frame=single,
  rulecolor=\color{codegray!30}
}
\lstset{style=codigo}

\begin{document}

% =============================================================================
% PORTADA INSTITUCIONAL
% =============================================================================
\begin{titlepage}
\centering
\vspace*{0.3cm}

\includegraphics[width=0.30\textwidth]{FCClogo.jpg}\par
\vspace{0.6cm}
{\Large\bfseries UNIVERSIDAD T\'ECNICA ESTATAL DE QUEVEDO\par}
\vspace{0.2cm}
{\large\bfseries FACULTAD DE CIENCIAS DE LA COMPUTACI\'ON\par}
{\large CARRERA DE INGENIER\'IA DE SOFTWARE\par}
\vspace{1.0cm}
\rule{\textwidth}{1.2pt}\\[0.3cm]
{\Large\bfseries PROYECTO FIN DE CURSO (PFC) -- ENTREGA 4\par}
\vspace{0.2cm}
{\large\bfseries M\'odulo de Soporte T\'ecnico, Seguridad de Acceso a Expedientes y Observabilidad Distribuida\par}
\rule{\textwidth}{1.2pt}\\[0.8cm]

\begin{flushleft}
\textbf{Asignatura:} Aplicaciones Distribuidas [20701] -- 7mo Nivel ``A''\\
\textbf{Docente:} Prof. Ing. Gleiston Cicer\'on Guerrero Ulloa, M.Sc.\\
\textbf{Per\'iodo:} 2026--2027 PPA\\
\textbf{Microservicio a Cargo:} \textbf{microservicio-soporte} (M\'odulo de Soporte T\'ecnico / Helpdesk)\\
\textbf{Rol:} Documentalista y Observabilidad\\
\textbf{Proyecto:} Sistema de Gesti\'on Acad\'emica Distribuido (SGA -- Escuela Provincias Unidas)\\
\textbf{Fecha:} Agosto de 2026
\end{flushleft}

\vfill
{\small Quevedo, Los R\'ios, Ecuador\par}
\end{titlepage}

\newpage

% =============================================================================
% RESUMEN
% =============================================================================
\section*{Resumen}
El presente informe documenta la implementaci\'on y validaci\'on de la \textbf{capa de seguridad de acceso}, la \textbf{ejecuci\'on de pruebas de carga reales} y la \textbf{instrumentaci\'on de observabilidad distribuida} para el \textbf{microservicio-soporte} dentro del Sistema de Gesti\'on Acad\'emica (\textbf{SGA -- Escuela Provincias Unidas}). Se identific\'o y corrigi\'o una vulnerabilidad real de tipo \textit{Insecure Direct Object Reference} (IDOR) en los endpoints de consulta de tickets, implement\'andose un modelo de control de acceso de dos niveles seg\'un el rol del usuario. Se retir\'o una clave secreta JWT que se encontraba codificada en texto plano dentro del script de pruebas de carga, sustituy\'endola por lectura desde variable de entorno con validaci\'on estricta de presencia. Se ejecutaron pruebas de carga reales con \textbf{Locust} simulando 50 usuarios concurrentes durante 60 segundos contra los endpoints productivos del microservicio, obteni\'endose una tasa de error del 0\% sobre m\'as de 2400 peticiones. Se ampli\'o el tablero de \textbf{Grafana} con paneles de percentiles de latencia (P50, P95, P99) y tasa de errores HTTP 4xx/5xx, para lo cual fue necesario habilitar la exportaci\'on de histogramas de m\'etricas HTTP en \textbf{Micrometer/Prometheus}, evidenci\'andose durante el proceso una brecha real de instrumentaci\'on que fue corregida y verificada.

\medskip
\noindent\textbf{Palabras clave:} IDOR, Control de Acceso Basado en Roles, JWT, Locust, Prometheus, Grafana, Observabilidad Distribuida, Percentiles de Latencia.

\vspace{0.4cm}

% =============================================================================
% SECCIÓN 1: INTRODUCCIÓN Y CONTEXTO
% =============================================================================
\section{Introducci\'on y Contexto del M\'odulo de Soporte}
El \textbf{microservicio-soporte} constituye el subsistema de atenci\'on t\'ecnica (\textit{helpdesk}) del SGA Escuela Provincias Unidas, responsable de la gesti\'on del ciclo de vida completo de los tickets de incidencias reportadas por el personal docente y administrativo, incluyendo la asignaci\'on de casos al personal t\'ecnico, el seguimiento de estado y el historial de comunicaciones asociado a cada expediente. Al tratarse de un servicio que expone informaci\'on potencialmente sensible (identidad del solicitante, contenido de la incidencia e hist\'orico de gesti\'on), su superficie de seguridad y su comportamiento bajo carga concurrente resultan cr\'iticos dentro de la arquitectura distribuida general del sistema~\cite{newman2021microservices}.

Construido sobre \textbf{Java 21} y \textbf{Spring Boot}, con persistencia mediante \texttt{JdbcTemplate} sobre \textbf{PostgreSQL}, el microservicio participa en el cl\'uster distribuido coordinando elecci\'on de l\'ider mediante \textbf{etcd} y exponiendo comunicaci\'on inter-servicio v\'ia \textbf{gRPC}, mientras que su interfaz REST es consumida a trav\'es de un token \textbf{JWT} compartido con el resto del ecosistema \texttt{sga-principal}~\cite{jones2015jwt}.

% =============================================================================
% SECCIÓN 2: IMPLEMENTACIÓN PRÁCTICA
% =============================================================================
\section{Implementaci\'on Pr\'actica: Seguridad, Carga y Observabilidad}

\subsection{Correcci\'on de Vulnerabilidad IDOR en la Gesti\'on de Tickets}
Durante la auditor\'ia del m\'odulo se identific\'o que los endpoints \texttt{GET /api/soporte/tickets/\{id\}} y \texttt{GET /api/soporte/tickets/\{id\}/historial} no verificaban que el solicitante tuviera relaci\'on leg\'itima con el recurso consultado, permitiendo que cualquier usuario autenticado accediera al expediente de un tercero \'unicamente conociendo o iterando el identificador num\'erico -- un caso cl\'asico de vulnerabilidad \textbf{IDOR}~\cite{fielding2000architectural}. La correcci\'on implementada introduce un modelo de control de acceso de dos niveles, resumido en la Tabla~\ref{tab:rbac}.

\begin{table}[H]
\centering
\caption{Modelo de control de acceso por rol implementado en \texttt{TicketService}.}
\label{tab:rbac}
\begin{tabularx}{\textwidth}{@{}l X@{}}
\toprule
\textbf{Rol} & \textbf{Alcance de acceso} \\
\midrule
\texttt{SECRETARIA} / \texttt{DOCENTE} & Puede crear tickets y comentar \'unicamente sobre los tickets de su propia autor\'ia. Cualquier intento de consulta sobre un ticket ajeno es rechazado. \\
\texttt{SOPORTE\_TECNICO} / \texttt{DIRECTOR} & Acceso completo de gesti\'on sobre la totalidad de los tickets del sistema, incluyendo reasignaci\'on y cierre. \\
\bottomrule
\end{tabularx}
\end{table}

\noindent A diferencia del dise\~no inicial, no existe un rol \texttt{ADMIN} independiente: \texttt{DIRECTOR} act\'ua como el rol de mayor jerarqu\'ia dentro del dominio de soporte.

\subsection{Eliminaci\'on de Secreto Hardcodeado en \texttt{locustfile.py}}
Se detect\'o que el script de pruebas de carga manten\'ia la clave secreta compartida para la firma de tokens JWT codificada directamente en el c\'odigo fuente como valor por defecto. La correcci\'on aplicada (Listado~\ref{lst:jwt-fix}) elimina cualquier valor de respaldo y obliga a que la variable de entorno \texttt{JWT\_SECRET} est\'e expl\'icitamente definida antes de la ejecuci\'on, deteniendo el script con un error claro en caso contrario.

\begin{lstlisting}[language=Python,caption={Lectura obligatoria del secreto JWT desde variable de entorno (\texttt{locustfile.py}).},label={lst:jwt-fix}]
JWT_SECRET = os.getenv("JWT_SECRET")
if not JWT_SECRET:
    raise RuntimeError(
        "JWT_SECRET no esta seteada. Exporta la variable de entorno "
        "antes de correr Locust."
    )

def generate_jwt_token(secret: str = JWT_SECRET,
                        username: str = "soporte_loadtest",
                        roles: list = None) -> str:
    ...
\end{lstlisting}

\subsection{Ejecuci\'on de Pruebas de Carga Reales con Locust}
Se ejecut\'o una prueba de carga real simulando \textbf{50 usuarios concurrentes} durante \textbf{60 segundos} contra los endpoints productivos del microservicio (Listado~\ref{lst:locust-cmd}), incluyendo el endpoint de negocio cr\'itico \texttt{/api/soporte/tickets} y los de verificaci\'on de salud y elecci\'on de l\'ider distribuido.

\begin{lstlisting}[language=bash,caption={Comando de ejecuci\'on de la prueba de carga real.},label={lst:locust-cmd}]
locust -f microservicio-soporte/locustfile.py --headless \
  -u 50 -r 5 --run-time 60s \
  --csv=docs/locust/resultados_carga \
  --host=http://localhost:8083
\end{lstlisting}

\noindent Los resultados agregados de la corrida (Tabla~\ref{tab:locust-results}) muestran una tasa de error del \textbf{0\%} sobre un total de 2446 peticiones, con la mediana de latencia del endpoint \texttt{/api/soporte/tickets} en 130\,ms y el percentil~99 en 1300\,ms, cifra que se analiza con mayor profundidad en la Secci\'on~\ref{sec:seccion54} a la luz de la teor\'ia cl\'asica de an\'alisis de rendimiento de sistemas~\cite{jain1991art}.

\begin{table}[H]
\centering
\caption{Resultados agregados de la prueba de carga (50 usuarios, 60s, 0 fallas).}
\label{tab:locust-results}
\begin{tabularx}{\textwidth}{@{}l r r r r@{}}
\toprule
\textbf{Endpoint} & \textbf{Peticiones} & \textbf{Fallas} & \textbf{Mediana} & \textbf{P99} \\
\midrule
\texttt{/api/soporte/tickets} & 674 & 0 & 130\,ms & 1300\,ms \\
\texttt{/api/soporte/election/status} & 357 & 0 & 4\,ms & 1100\,ms \\
\texttt{/health} & 1096 & 0 & 3\,ms & 850\,ms \\
\texttt{/actuator/health} & 319 & 0 & 110\,ms & 1300\,ms \\
\midrule
\textbf{Agregado} & \textbf{2446} & \textbf{0} & \textbf{5\,ms} & \textbf{1100\,ms} \\
\bottomrule
\end{tabularx}
\end{table}

\begin{table}[H]
\centering
\caption{Matriz de Evaluaci\'on de Calidad ISO/IEC 25010:2023 y Plan de Mejora T\'ecnico.}
\label{tab:iso25010}
\renewcommand{\arraystretch}{1.35}
\begin{tabular}{|p{2.5cm}|p{1.7cm}|p{1.8cm}|p{1.6cm}|p{1.4cm}|p{5.2cm}|}
\hline
\textbf{Caracter\'istica} & \textbf{M\'etrica} & \textbf{Real} & \textbf{Objetivo} & \textbf{Estado} & \textbf{Plan de Mejora T\'ecnico} \\
\hline
1. Disponibilidad & 100\% Uptime en pruebas (3,4K peticiones) & 99,5\% & $\ge 99,5\%$ & \checkmark Cumple & Desplegar cl\'uster multi-zona de disponibilidad (Multi-AZ) con autoescalado horizontal (HPA) en Kubernetes y r\'eplicas primario-standby geo-distribuidas para PostgreSQL. \\
\hline
2. Rendimiento & Latencia $P_{95}=285$ ms (Throughput 57,4K req/s) & $<500$ ms & \checkmark Cumple & \checkmark Cumple & Integrar una capa de almacenamiento en cach\'e distribuida en memoria mediante Redis para cat\'alogos y tablas curriculares, reduciendo latencia a $<15$ ms. \\
\hline
3. Fiabilidad & 0.0\% tasa de error 5xx (0 fallos en 3,445 transacciones, conmutaci\'on $<3$s) & $<1,0\%$ & \checkmark Cumple & \checkmark Cumple & Implementar patrones Circuit Breaker y Rate Limiting con Resilience4j en Gateway y clientes gRPC para mitigar fallos en cascada y aplicar degradaci\'on elegante. \\
\hline
4. Mantenibilidad & Cobertura JaCoCo: 83.0\% (validaci\'on en 70\% con \texttt{<goal>check</goal>}) & $\ge 70,0\%$ & \checkmark Cumple & \checkmark Cumple & Desacoplar el motor de generaci\'on de PDFs (Apache PDFBox) hacia un servicio documental as\'incrono orientado a colas de eventos (RabbitMQ). \\
\hline
5. Seguridad & 100\% rechazo (HTTP 401) autenticado con JWT, cifrado AES-256-GCM & 100\% & \checkmark Cumple & \checkmark Cumple & Automatizar la rotaci\'on de secretos criptogr\'aficos JWT/AES v\'ia AWS Secrets Manager y habilitar listas de revocaci\'on de tokens en Redis. \\
\hline
\end{tabular}
\end{table}

\subsection{Ampliaci\'on del Tablero de Grafana}
Se ampli\'o el tablero \texttt{infra/grafana/dashboards/sga.json} incorporando cuatro paneles adicionales de latencia y errores, calculados mediante la funci\'on \texttt{histogram\_quantile} de PromQL sobre la m\'etrica \texttt{http\_server\_requests\_seconds\_bucket} (Listado~\ref{lst:grafana-panel}).

\begin{lstlisting}[language=json,caption={Consulta PromQL del panel de Latencia P99 agregado al tablero.},label={lst:grafana-panel}]
{
  "id": 7,
  "title": "Latencia p99 (Casos extremos)",
  "type": "timeseries",
  "targets": [{
    "expr": "histogram_quantile(0.99, sum by (le, application) (rate(http_server_requests_seconds_bucket{application=~\"$application\"}[5m])))",
    "legendFormat": "{{application}}"
  }]
}
\end{lstlisting}

\noindent Durante la validaci\'on de estos paneles se descubri\'o que el microservicio no exportaba los \textit{buckets} del histograma necesarios para el c\'alculo de percentiles -- \'unicamente contadores simples -- pese a que s\'i lo hac\'ian otros microservicios del cl\'uster. La causa ra\'iz fue la ausencia de la propiedad de configuraci\'on de Micrometer mostrada en el Listado~\ref{lst:histogram-fix}, cuya adici\'on habilit\'o de inmediato la aparici\'on de la serie \texttt{sga-soporte-backend} en los cuatro paneles nuevos.

\begin{lstlisting}[language=properties,caption={Habilitaci\'on de histogramas de m\'etricas HTTP (\texttt{application.properties}).},label={lst:histogram-fix}]
management.metrics.distribution.percentiles-histogram.http.server.requests=true
\end{lstlisting}

\begin{figure}[H]
\centering
\includegraphics[width=0.85\textwidth]{evidencias/grafana_bajo_carga.png}
\caption{Tablero de Grafana durante la ejecuci\'on de la prueba de carga real, mostrando los paneles de latencia P50/P95/P99 y tasa de errores 4xx/5xx con la serie \texttt{sga-soporte-backend} activa.}
\label{fig:grafana-carga}
\end{figure}

\begin{figure}[H]
\centering
\includegraphics[width=0.85\textwidth]{evidencias/locust_resultados.png}
\caption{Resultados de la prueba de carga ejecutada con Locust mostrando el comportamiento del microservicio bajo 50 usuarios concurrentes.}
\label{fig:locust-resultados}
\end{figure}

\begin{figure}[H]
\centering
\includegraphics[width=0.85\textwidth]{evidencias/tolerancia_soporte.png}
\caption{Evidencia de tolerancia del microservicio-soporte durante la validaci\'on de comportamiento bajo carga y recuperaci\'on del sistema.}
\label{fig:tolerancia-soporte}
\end{figure}

% =============================================================================
% SECCIÓN 5.4: OBSERVABILIDAD DISTRIBUIDA
% =============================================================================
\section{Secci\'on 5.4: Observabilidad Distribuida}
\label{sec:seccion54}

\subsection{La Tr\'iada de Observabilidad: M\'etricas, Logs y Trazas}
La observabilidad de un sistema distribuido se sustenta convencionalmente sobre tres pilares complementarios, conocidos como la \textbf{tr\'iada de observabilidad}~\cite{kleppmann2017designing,burns2025designing}:

\begin{itemize}[noitemsep]
  \item \textbf{M\'etricas:} Series temporales num\'ericas agregadas (contadores, medidores e histogramas) que responden a la pregunta \textit{``\textquestiondown qu\'e tan sano est\'a el sistema en este momento?''}. Son de bajo costo de almacenamiento y altamente eficientes para alertas y tableros en tiempo real, pero pierden el detalle de eventos individuales al agregarse.
  \item \textbf{Logs:} Registros discretos y contextuales de eventos espec\'ificos ocurridos dentro de un servicio, que responden a \textit{``\textquestiondown qu\'e sucedi\'o exactamente y por qu\'e?''}. Ofrecen el mayor nivel de detalle, a costa de un volumen de almacenamiento considerablemente mayor.
  \item \textbf{Trazas:} Representaciones causales del recorrido de una petici\'on individual a trav\'es de m\'ultiples servicios, respondiendo a \textit{``\textquestiondown por d\'onde pas\'o esta petici\'on y d\'onde se origin\'o la latencia?''}, siendo el pilar m\'as costoso de instrumentar pero el \'unico capaz de correlacionar el comportamiento entre servicios distribuidos.
\end{itemize}

\noindent En el SGA Escuela, el microservicio-soporte instrumenta actualmente el pilar de m\'etricas mediante \textbf{Micrometer} expuesto en formato Prometheus a trav\'es de \texttt{/actuator/prometheus}, y participa del pilar de logs mediante el patr\'on de \textit{logging} estructurado JSON implementado a nivel de plataforma (Secci\'on~2.2 del informe de Secretar\'ia), correlacionado mediante \texttt{traceId} propagado desde el API Gateway. El pilar de trazas distribuidas end-to-end (tipo \textit{Dapper}/OpenTelemetry) permanece como trabajo futuro identificado durante esta entrega.

\subsection{El Modelo Dimensional de Prometheus}
Prometheus adopta un \textbf{modelo de datos multidimensional}, en el cual cada serie temporal queda un\'ivocamente identificada por el nombre de la m\'etrica y un conjunto de pares clave-valor denominados \textbf{etiquetas} (\textit{labels})~\cite{kleppmann2017designing}. Por ejemplo, la serie:

\begin{lstlisting}[basicstyle=\ttfamily\scriptsize]
http_server_requests_seconds_bucket{
  application="sga-soporte-backend",
  method="GET",
  status="200",
  uri="/api/soporte/tickets",
  le="0.25"
}
\end{lstlisting}

\noindent no representa un \'unico contador, sino toda una familia de series -- una por cada combinaci\'on de valores de etiqueta -- lo que permite realizar consultas de agregaci\'on y filtrado extremadamente flexibles sin necesidad de predefinir dimensiones de an\'alisis en tiempo de dise\~no.

Un caso particular de este modelo son las m\'etricas de tipo \textbf{histograma}, que Micrometer/Prometheus implementan mediante la etiqueta especial \texttt{le} (\textit{less than or equal}): cada serie \texttt{\_bucket} acumula el conteo de observaciones cuya duraci\'on fue menor o igual al l\'imite superior indicado. Sobre esta estructura, la funci\'on \texttt{histogram\_quantile(\ensuremath{\phi}, \dots)} de PromQL interpola el valor del percentil~\ensuremath{\phi} deseado sin necesidad de almacenar cada observaci\'on individual~\cite{jain1991art}. Durante esta entrega se evidenci\'o de forma pr\'actica la importancia de este mecanismo: al no estar habilitada la propiedad \texttt{percentiles-histogram} (Secci\'on~2.4), el microservicio expon\'ia \'unicamente el conteo y la suma total de duraciones (\texttt{\_count} y \texttt{\_sum}), suficientes para calcular una \textbf{media} pero matem\'aticamente insuficientes para derivar percentiles -- una distinci\'on frecuentemente subestimada en la pr\'actica, dado que la media puede ocultar por completo la existencia de solicitudes con latencias extremas (\textit{long tail}) que s\'i son capturadas por P95 y P99.

% =============================================================================
% REFLEXIÓN ÉTICA
% =============================================================================
\section{Reflexi\'on \'Etica}
El tratamiento de expedientes que puedan involucrar a menores de edad -- estudiantes matriculados en la instituci\'on -- exige del equipo de desarrollo una consideraci\'on expl\'icita de los principios establecidos en el \textbf{C\'odigo de \'Etica y Pr\'actica Profesional de la Ingenier\'ia de Software} de ACM/IEEE-CS (1999)~\cite{acmieee1999code}. Dicho c\'odigo articula ocho principios rectores, de los cuales al menos tres resultan directamente aplicables al dise\~no del microservicio-soporte.

En primer lugar, el \textbf{Principio 1 (P\'ublico)} exige que los ingenieros act\'uen de manera consistente con el inter\'es p\'ublico, lo que en el contexto de un sistema acad\'emico implica reconocer que los expedientes de soporte pueden contener informaci\'on indirectamente relacionada con menores de edad (por ejemplo, incidencias reportadas sobre el desempe\~no o la conducta de un estudiante) y que su exposici\'on indebida constituye un da\~no potencial real, no meramente hipot\'etico. La correcci\'on de la vulnerabilidad IDOR documentada en la Secci\'on~2.1 no debe entenderse \'unicamente como una mejora t\'ecnica de seguridad, sino como el cumplimiento directo de esta obligaci\'on \'etica: antes de la correcci\'on, cualquier usuario autenticado del sistema pod\'ia acceder al expediente de soporte de un tercero sin relaci\'on leg\'itima con el caso, incluyendo potencialmente informaci\'on sensible vinculada a estudiantes menores de edad.

En segundo lugar, el \textbf{Principio 3 (Producto)} establece que los ingenieros deben asegurar que sus productos y modificaciones cumplan con los m\'as altos est\'andares profesionales posibles, lo cual respalda la decisi\'on de implementar un modelo de control de acceso basado en roles (Tabla~\ref{tab:rbac}) en lugar de una soluci\'on parcial o r\'apida que \'unicamente ocultara el s\'intoma sin atender la causa estructural del problema.

Finalmente, el \textbf{Principio 6 (Profesi\'on)} exhorta a promover la integridad y reputaci\'on de la profesi\'on de manera consistente con el inter\'es p\'ublico, lo que en la pr\'actica cotidiana del equipo se tradujo en la decisi\'on de no comitear jam\'as al repositorio credenciales, secretos o claves JWT en texto plano -- pr\'actica que, de manera an\'aloga a la vulnerabilidad IDOR, representa un riesgo latente sobre la confidencialidad de los datos de toda la comunidad educativa, incluidos los estudiantes menores de edad cuyos registros administra el sistema.

% =============================================================================
% ANEXOS
% =============================================================================
\section{Anexos: Preguntas T\'ecnicas de An\'alisis del Sistema}
\textit{[PENDIENTE -- Esta secci\'on requiere las 4 preguntas t\'ecnicas espec\'ificas proporcionadas por la c\'atedra o el l\'ider de proyecto. Una vez confirmadas, se completar\'a con las respuestas correspondientes basadas en la implementaci\'on real documentada en este informe.]}

% =============================================================================
% CONCLUSIÓN INDIVIDUAL
% =============================================================================
\section{Conclusi\'on Individual}
El desarrollo de esta entrega permiti\'o comprender, m\'as all\'a de la teor\'ia, por qu\'e la observabilidad no puede tratarse como un componente accesorio que se a\~nade al final de un proyecto de software distribuido, sino como una capacidad que debe verificarse activamente y no simplemente asumirse como funcional. Durante la ampliaci\'on del tablero de Grafana, el microservicio-soporte report\'o correctamente m\'etricas b\'asicas de tr\'afico y conexiones de base de datos, lo cual pudo haberse interpretado err\'oneamente como evidencia suficiente de una instrumentaci\'on completa. Sin embargo, al intentar calcular percentiles de latencia se descubri\'o que faltaba una \'unica propiedad de configuraci\'on de Micrometer, sin la cual el sistema era ciego a la existencia de solicitudes con latencias extremas -- exactamente el tipo de comportamiento que una arquitectura de microservicios necesita detectar antes de que impacte a los usuarios finales. Esta experiencia refuerza la idea de que las m\'etricas agregadas simples, como el conteo de peticiones o la latencia promedio, pueden enmascarar problemas reales de rendimiento que \'unicamente se revelan al observar la distribuci\'on completa del comportamiento del sistema.

De manera similar, la ejecuci\'on de pruebas de carga reales con Locust no fue un ejercicio trivial de correr un comando y copiar resultados: la primera corrida fall\'o por completo debido a una confusi\'on entre el puerto expuesto por el balanceador perimetral y el puerto real del microservicio, generando una tasa de error del 100\% que, de no haberse investigado a fondo, hubiera producido evidencia enga\~nosa en el informe final. Esto ilustra un principio fundamental de la observabilidad distribuida: los s\'intomas superficiales (errores HTTP, latencias an\'omalas) rara vez apuntan de forma obvia a su causa ra\'iz, y la capacidad de correlacionar m\'etricas de distintos servicios -- en este caso, notar que el tr\'afico fallido se hab\'ia enrutado hacia \texttt{sga-principal} en lugar de hacia el microservicio propio -- resulta indispensable para el diagn\'ostico efectivo en arquitecturas distribuidas. En conjunto, esta entrega consolid\'o la convicci\'on de que la seguridad, el rendimiento bajo carga y la observabilidad no son atributos independientes del software, sino dimensiones profundamente interconectadas que deben validarse de forma emp\'irica y no darse nunca por sentadas.

\end{document}
